\paragraph{Short-horizon trajectory prediction.}
Most trajectory prediction work targets short horizons: given a brief observation window, predict near-future displacements.
Modern architectures (e.g., Transformers \citep{vaswani2017attention}) and generative variants have pushed accuracy forward, but these settings typically do not require committing to a full-trip route under origin--destination intent.
In contrast, our focus is full-trip route generation, where corridor-level alternatives and feasibility constraints dominate the difficulty.

\paragraph{End-to-end generative models for long-horizon routes.}
A natural extension is to train an end-to-end generator to output an entire coordinate sequence, using latent-variable models \citep{kingma2014auto} or diffusion models \citep{ho2020ddpm}.
However, in route generation the output distribution is strongly multi-modal in a \emph{topological} sense: distinct corridor-level modes can be geometrically far apart yet semantically equivalent.
This makes end-to-end coordinate generation prone to destructive averaging and mode interference, which we highlight via diagnostic experiments.

\paragraph{Hierarchical and latent-structure approaches.}
Hierarchical approaches address long-horizon uncertainty by separating coarse decisions from fine execution, e.g., sampling sparse waypoints/anchors and then refining a continuous path.
Our formulation instantiates this idea as an explicit decision--execution cascade: a discrete topological commitment (waypoint skeleton) precedes continuous route generation, so corridor-level multi-modality is represented explicitly rather than being forced into a single continuous output space.
Beyond introducing the hierarchy, we emphasize a composition issue that is often glossed over in practice: a powerful continuous refiner can overwrite the coarse commitment, collapsing corridor diversity even when the decision stage is multi-modal.

\paragraph{Feasibility: hard constraints vs soft priors.}
Feasibility is often enforced by hard conditioning on a road graph or by masking the output space to a ``feasible set''.
Hard constraints can ensure on-road outputs but may also truncate the learned distribution and bind performance to map quality.
We instead use map information as soft priors (e.g., road proximity fields) to encourage plausibility without hard support cutting.

\paragraph{Macro priors and population synthesis (related but distinct).}
Macro-level statistics (e.g., census attributes) are central in population synthesis and controllable urban simulation.
Our work touches this direction only lightly: census-derived spatial covariates are explored as an optional guidance signal to demonstrate extensibility, while the core contribution remains the generative architecture for route generation.

\paragraph{Note on citations (draft).}
This draft focuses on clarifying the core claims and diagnostic evidence; we will expand citations to domain-specific route choice, trajectory generation, and hierarchical planning work in the final version.
